\section{Computing expected values}

\begin{lem}
If $X$ is a discrete r.v., that is, takes values in
some countable set $S$, then

\[
\mathbf{E} X = \sum_{s\in S} s\, \mathrm{P}(X=s)
\]

\noindent when the right-hand side is well-defined, i.e., when
at least one of the numbers

\[
\mathbf{E}(X_-)=\sum_{s\in S, s<0} (-s)\,\mathrm{P}(X=s), \qquad \mathbf{E}(X_+)= \sum_{s\in S, s>0} s\, \mathrm{P}(X=s)
\]

\noindent is finite. It follows that for any function
$g:\R\to\R$, we also have

\[
\mathbf{E}(g(X)) = \sum_{s\in S} g(s) \mathrm{P}(X=s).
\]

\end{lem}

\begin{proof}
If $S$ is finite then $X$ is a simple function, and
can be written $X=\sum_{s\in S} s\,1_{\{X=s\}}$, so
this follows from the definition of
$\mathbf{E}(\cdot)$ for simple functions. If $S$ is
infinite this follows (check!) from the convergence
theorems in the previous section by considering
approximations to $X$ of the form
$\sum_{s\in S, |s|<M} s\, 1_{\{X=s\}}$.
\end{proof}

\begin{lem}
If $X$ is a r.v. with a density function $f_X$, then

\[
\mathbf{E}(X) = \int_{-\infty}^\infty x f_X(x)\,dx
\]

\noindent when the right-hand side is well-defined, i.e., when
at least one of the numbers

\[
\mathbf{E}(X_-) = -\int_{-\infty}^0 x f_X(x)\,dx, \qquad \mathbf{E}(X_+) = \int_0^\infty x f_X(x)\,dx
\]

\noindent is finite. Similarly, for any ``reasonable'' function
$g$ we have

\[
\mathbf{E}(g(X)) = \int_{-\infty}^\infty g(x) f_X(x)\, dx.
\]

\end{lem}

\begin{proof}
Fix $\epsilon>0$, and approximate $X$ by a discrete
r.v. $Y$, e.g.,

\[
Y = \sum_{k=-\infty}^\infty k\epsilon 1_{\{k\epsilon < X \le (k+1)\epsilon\}}.
\]

\noindent Then
$|\mathbf{E}(X)-\mathbf{E}(Y)| \le \mathbf{E}|X-Y|
\le \epsilon$.
By the previous lemma we have

\[
\mathbf{E}(Y)  = \sum_{k=-\infty}^\infty k\epsilon \mathrm{P}(k\epsilon<X<(k+1)\epsilon) =\sum_{k=-\infty}^\infty k\epsilon \int_{k\epsilon}^{(k+1)\epsilon} f_X(x)\,dx,
\]

\noindent so the result for $\mathbf{E} X$ follows by letting
$\epsilon\to 0$. For general functions $g(X)$ repeat
this argument, and invoke the relevant convergence
theorem to deduce that
$\mathbf{E}(g(Y))\to \mathbf{E} g(Y)$ as
$\epsilon\to 0$.
\end{proof}
